\documentclass[10pt,a4paper]{amsart}
\usepackage[margin=19mm]{geometry}
\usepackage{amsmath,amssymb,mathtools}
\usepackage[scr=rsfs,cal=euler]{mathalfa}
\usepackage{xfrac}
\usepackage[unicode,hidelinks,bookmarks=false]{hyperref}
\newcommand{\ie}{i.e.\ }
\newcommand{\cf}{cf.\ }
\newcommand{\resp}{respectively\ }
\newcommand{\Subsection}{\subsection}
\providecommand{\nobreakdash}{\nobreak\hbox{-}}
\setlength{\emergencystretch}{2em}
\multlinegap0pt
\usepackage{amssymb}
\usepackage{xcolor}
\usepackage[shortlabels]{enumitem}
\usepackage{float}
\newtheorem{lemma}{Lemma}
\newtheorem{proposition}{Proposition}
\newtheorem{theorem}{Theorem}
\newcommand{\F}{\mathbb{F}}
\newcommand{\abs}[1]{\left|#1\right|}
\title{Parabolic distance in $\F_q^2$: a sharp exponent and new results}
\author{Dao Nguyen Van Anh, Steven Senger, Dung The Tran, Le Anh Vinh}
\date{Source: arXiv:2603.21292v1 (CC BY 4.0)}
\begin{document}
\maketitle

\noindent\emph{Source and licence.} Parabolic distance in $\F_q^2$: a sharp exponent and new results. Version arXiv:2603.21292v1 (CC BY 4.0), distributed by arXiv under the Creative Commons Attribution 4.0 International licence, http://creativecommons.org/licenses/by/4.0/, as recorded in the arXiv licence metadata for this exact version and shown on the abstract page https://arxiv.org/abs/2603.21292v1. Full licence notice: \texttt{licenses/arxiv-2603.21292-CC-BY-4.0.txt}.\par\smallskip

\noindent\textit{Editorial scope.} Counted unit: the article's theorem thm:bipartite (source0.tex lines 772--796). The article's complete original proof of it is delivered with it (source0.tex lines 798--969, one \texttt{proof} environment of 172 lines). No mathematics is written, repaired or re-derived: the statement and the proof are the article's own lines, in the article's own notation, and no retained line is substituted or re-flowed.  Retained uncounted as the source-local context the proof needs: lines 143--145; lines 174--180; lines 189--193; lines 202--211; lines 487--491; lines 538--545; lines 558--562; lines 564--568.  Nothing was omitted from any retained span, so the package carries no omission record.  No retained line contains a citation, so the wrapper carries no bibliography and no bibliography entry.  No retained line contains a figure environment, an image-inclusion command or drawing code, so no image is delivered and the package declares no asset dependency.  Editorial formatting only, in addition: an AMS \texttt{amsart} preamble that restates the article's own declarations and macros used by the retained lines, namely the environments \texttt{lemma}, \texttt{proposition}, \texttt{theorem} on separate counters, together with the standard packages those lines need.  The retained environments print in this document as Lemma 1, Proposition 1, Theorem 1 in order of appearance, whereas the article prints the same items with the numbering of its own full document.  The complete native source of the article as distributed in the arXiv e-print for this exact version is bundled unchanged as \texttt{sources/196935/source0.tex}.
\begin{equation}\label{eq:def-K}
K_E:=\max_{u\in \F_q}\abs{E\cap (\{u\}\times \F_q)},
\end{equation}

\begin{equation}\label{eq:orthogonality}
\sum_{t\in \F_q}\chi(\alpha t)=
\begin{cases}
q, & \alpha=0,\\
0, & \alpha\ne 0.
\end{cases}
\end{equation}

\begin{equation}\label{lem:plancherel}
\sum_{\xi\in \F_q}\abs{\widehat{h}(\xi)}^2
=
q\sum_{y\in \F_q}\abs{h(y)}^2.
\end{equation}

\begin{lemma}\label{lem:energy-trivial}
For finite sets $P,Q\subset \F_q$,
\begin{equation*}
E_+(P,Q)
\le
|P|\,|Q|\,\min\{|P|,|Q|\}
\le
|P|^{\frac{3}{2}}|Q|^{\frac{3}{2}}.
\end{equation*}
\end{lemma}

\begin{equation}\label{eq:def-shears}
A:=\{(x_1,x_2+x_1^2):\mathbf{x}=(x_1,x_2)\in E\},
\qquad
B:=\{(y_1,y_2-y_1^2):\mathbf{y}=(y_1,y_2)\in E\}.
\end{equation}

\begin{proposition}\label{prop:second-moment}
With $n:=|E|$ and $K_E$ as in \eqref{eq:def-K},
\begin{equation}\label{eq:second-moment-bound}
\sum_{t\in \F_q}\nu(t)^2
\le
\frac{n^4}{q}+qn^2K_E.
\end{equation}
\end{proposition}

\begin{equation}\label{eq:nu-split}
\sum_{t\in \F_q}\nu(t)^2
=
\frac{n^4}{q}+\sum_{t\in \F_q}R_t^2.
\end{equation}

\begin{equation}\label{eq:R-by-S}
\sum_{t\in \F_q}R_t^2
=
\frac{1}{q}\sum_{s\ne 0}|S(s)|^2.
\end{equation}

\begin{theorem}\label{thm:bipartite}
Let $E,F\subset \F_q^2$. Then
\[
\sum_{t\in\F_q}\nu_{E,F}(t)^2
\le
\frac{|E|^2|F|^2}{q}
+
q\,|E|\,|F|\,\sqrt{K_EK_F},
\]
where
\[
\nu_{E,F}(t):=
\abs{\{(\mathbf{x},\mathbf{y})\in E\times F:\ \|\mathbf{x}-\mathbf{y}\|_P=t\}}.
\]
Consequently,
\[
\abs{\Delta_P(E,F)}
\ge
\frac{|E|^2|F|^2}
{\dfrac{|E|^2|F|^2}{q}+q\,|E|\,|F|\,\sqrt{K_EK_F}}
=
\frac{q}{1+\dfrac{q^2\sqrt{K_EK_F}}{|E|\,|F|}}.
\]
In particular, if $|E|\,|F|\gg q^2\sqrt{K_EK_F}$, then $\abs{\Delta_P(E,F)}\gg q$.
\end{theorem}

\begin{proof}
We follow the proof of Proposition~\ref{prop:second-moment}, replacing the two copies of $E$ by the pair $E,F$.

For $t\in \F_q$, define
\[
\nu_{E,F}(t):=
\abs{\{(\mathbf{x},\mathbf{y})\in E\times F:\ \|\mathbf{x}-\mathbf{y}\|_P=t\}}.
\]
Then
\[
\sum_{t\in\F_q}\nu_{E,F}(t)=|E|\,|F|.
\]
Hence, by Cauchy--Schwarz,
\begin{equation}\label{eq:bipartite-CS}
|E|^2|F|^2
=
\left(\sum_{t\in\F_q}\nu_{E,F}(t)\right)^2
\le
\abs{\Delta_P(E,F)}
\sum_{t\in\F_q}\nu_{E,F}(t)^2,
\end{equation}
so it remains to bound the second moment $\sum\limits_{t\in\F_q}\nu_{E,F}(t)^2$.

As in \eqref{eq:def-shears}, define
\[
A:=\{(x_1,x_2+x_1^2):\ \mathbf{x}=(x_1,x_2)\in E\},
\qquad
B:=\{(y_1,y_2-y_1^2):\ \mathbf{y}=(y_1,y_2)\in F\}.
\]
Then $|A|=|E|$ and $|B|=|F|$, and for $\mathbf{x}\in E$, $\mathbf{y}\in F$, with corresponding points
$\mathbf{a}\in A$, $\mathbf{b}\in B$, one has
\[
\|\mathbf{x}-\mathbf{y}\|_P
=
(x_2-y_2)+(x_1-y_1)^2
=
a_2-b_2-2a_1b_1.
\]
Therefore
\[
\nu_{E,F}(t)
=
\abs{\{(\mathbf{a},\mathbf{b})\in A\times B:\ a_2-b_2-2a_1b_1=t\}}.
\]

Using the orthogonality relation \eqref{eq:orthogonality}, we obtain
\[
\nu_{E,F}(t)
=
\frac{1}{q}\sum_{s\in \F_q}\chi(-st)S(s),
\]
where
\[
S(s):=
\sum_{\mathbf{a}\in A}\sum_{\mathbf{b}\in B}
\chi\bigl(s(a_2-b_2-2a_1b_1)\bigr).
\]
The $s=0$ term contributes $\frac{|E| |F|}{q}$. Setting
\[
R_t:=\nu_{E,F}(t)-\frac{|E| |F|}{q},
\]
the same computation as in \eqref{eq:nu-split} and \eqref{eq:R-by-S} gives
\[
\sum_{t\in\F_q}\nu_{E,F}(t)^2
=
\frac{|E|^2|F|^2}{q}
+
\frac{1}{q}\sum_{s\ne 0}|S(s)|^2.
\]
Thus it remains to estimate $\sum\limits_{s\ne 0}|S(s)|^2$.

For $x,y\in\F_q$, define the vertical slices
\[
A_x:=\{u\in\F_q:(x,u)\in A\},
\qquad
B_y:=\{v\in\F_q:(y,v)\in B\},
\]
and for $s\ne 0$ set
\[
f_s(x):=\sum_{u\in A_x}\chi(su),
\qquad
g_s(y):=\sum_{v\in B_y}\chi(-sv).
\]
Then
\[
S(s)=\sum_{x\in\F_q}f_s(x)\,\widehat{g_s}(-2sx).
\]
Define
\[
U_s:=\sum_{x\in\F_q}|f_s(x)|^2,
\qquad
V_s:=\sum_{y\in\F_q}|g_s(y)|^2.
\]
Exactly as in the proof of Proposition~\ref{prop:second-moment}, the Cauchy--Schwarz inequality and Plancherel \eqref{lem:plancherel} yield
\[
|S(s)|^2\le q\,U_sV_s
\qquad (s\ne 0),
\]
and hence
\begin{equation}\label{eq:bipartite-sumS}
\sum_{s\ne 0}|S(s)|^2
\le
q\left(\sum_{s\in\F_q}U_s^2\right)^{\frac{1}{2}}
\left(\sum_{s\in\F_q}V_s^2\right)^{\frac{1}{2}}.
\end{equation}

We now estimate the two factors separately. By the same expansion used in the proof of
Proposition~\ref{prop:second-moment},
\[
\sum_{s\in\F_q}U_s^2
=
q\sum_{x,x'\in\F_q}E_+(A_x,A_{x'}).
\]
Applying Lemma~\ref{lem:energy-trivial},
\[
E_+(A_x,A_{x'})
\le
|A_x|^{\frac32}|A_{x'}|^{\frac32},
\]
so
\[
\sum_{s\in\F_q}U_s^2
\le
q\left(\sum_{x\in\F_q}|A_x|^{\frac32}\right)^2.
\]
Since the shear preserves vertical fibers, we have
\[
\max_{x\in\F_q}|A_x|=K_E
\qquad\text{and}\qquad
\sum_{x\in\F_q}|A_x|=|E|.
\]
Therefore
\[
\sum_{x\in\F_q}|A_x|^{\frac32}
\le
K_E^{\frac{1}{2}}\sum_{x\in\F_q}|A_x|
=
K_E^{\frac{1}{2}}|E|,
\]
and hence
\[
\sum_{s\in\F_q}U_s^2\le qK_E|E|^2.
\]
By the same argument applied to $B$,
\[
\sum_{s\in\F_q}V_s^2\le qK_F|F|^2.
\]
Substituting these bounds into \eqref{eq:bipartite-sumS}, we find
\[
\sum_{s\ne 0}|S(s)|^2
\le
q^2|E|\,|F|\,\sqrt{K_EK_F}.
\]
Consequently,
\[
\sum_{t\in\F_q}\nu_{E,F}(t)^2
\le
\frac{|E|^2|F|^2}{q}
+
q\,|E|\,|F|\,\sqrt{K_EK_F}.
\]
Finally, combining this with \eqref{eq:bipartite-CS} gives
\[
\abs{\Delta_P(E,F)}
\ge
\frac{|E|^2|F|^2}
{\dfrac{|E|^2|F|^2}{q}+q\,|E|\,|F|\,\sqrt{K_EK_F}}
=
\frac{q}{1+\dfrac{q^2\sqrt{K_EK_F}}{|E|\,|F|}},
\]
as claimed.
\end{proof}

\end{document}
